\abstract{
Ill-defined urban planning problems, such as master plan generation, involve conflicting objectives, heterogeneous data sources, ambiguous evaluation criteria, and no single correct solution.
Existing LLM-based multi-agent systems (LLM-MAS) are evaluated primarily on well-defined benchmarks with ground-truth answers, creating a fundamental mismatch between system capabilities and real-world decision-making needs.
We introduce FP2MP (From Prompt 2 Master Plan), a blackboard-based multi-agent architecture that operationalizes ill-defined problem-solving as iterative uncertainty reduction through a closed-loop workflow: decomposition (ReDI), execution via specialized tool-augmented agents (web search, normative retrieval, geospatial computation), curation with evidence promotion, and critic-driven follow-up injection.
The system maintains a three-level shared memory (raw data $\rightarrow$ wiki pages $\rightarrow$ confirmed facts), enabling adaptive collaboration depth without fixed interaction rounds.
We also propose FP2MP-Eval, a multi-dimensional LLM-as-a-judge evaluation protocol assessing problem structuring, solution quality, and reflective reasoning across eight criteria (5-point Likert).
Validation on 30 ill-defined urban planning tasks demonstrates high inter-rater reliability (ICC > 0.82) and sensitivity to model capability differences.
Quantitative comparison against six baselines (Single-Agent, Chain-of-Thought, ReAct, Actor-Critic, Majority Voting, Multi-Agent Debate) on 60 problem instances shows that FP2MP achieves superior performance in metacognition (5.00 $\pm$ 0.00 vs. 2.26–3.51 for baselines) and uncertainty handling (4.74 $\pm$ 0.45 vs. 2.46–3.83), while remaining competitive in coherence and decomposition.
A case study on Gatchina, Russia, demonstrates that FP2MP proposes a polycentric, evidence-based master plan grounded in verifiable spatial computations (e.g., clinic provision improvement from 0.295 to 0.833).
The open-source implementation supports reproducibility and provides a foundation for LLM-MAS evaluation on ill-defined problems beyond urban planning.
}